%=======================================================================
% Titelseite
%=======================================================================
\chapter{Theoretische Grundlagen}
\label{sec:BasicRulez}

\begin{figure}[h]
	\centering
		\includegraphics[width=0.80\textwidth]{Bilder/Python_God.jpg}
	\caption{Es gibt Regeln an die man sich halten sollte \cite{Python83}}
	\label{fig:Python_God}
\end{figure}
In diesem Unterkapitel "`Theoretische Grundlagen"' soll ein Überblick über die unterschiedlichen Definitionen und Ansätze zum jeweiligen Thema und eine Abgrenzung zu verwandten Themen und Begriffen erfolgen \cite[S.~27]{Karm2011}.
Wissenschaftliche Arbeiten sollen Klarheit beim Leser hervorrufen.
Das bedeutet, dass der Leser die beschriebenen Sachverhalte exakt in der beabsichtigten Weise ohne Rückfragen verstehen muss.
Welche Definitionen sind relevant um den Bericht ohne weitere Hintergrundkenntnisse in dem Themenbereich zu verstehen?
Inhalte müssen logisch gegliedert und klar nachvollziehbar sein.
Es ist ein hohes Maß an Systematik, Ordnung, Logik und Klarheit erforderlich \cite[S.~1~ff.]{Hering2009}.

"`Zwischendurch müssen Sie immer wieder einmal die gedankliche Position des Lesers einnehmen und sich fragen:
\begin{itemize}
	\item Wann braucht der Leser welche Information?
	\item Fällt das hier verwendete Bild 'vom Himmel'?
	\item Soll ich mich auf die Gliederung beziehen, eine Zwischenzusammenfassung machen oder den neuen Dokumentteil aus übergeordneter Sicht ankündigen?
	\item Ist die Aufteilung der Informationen logisch und nachvollziehbar?"' \cite[S.~29]{Hering2009}.
\end{itemize}

Die nachfolgenden Grundregeln sind strikt zu beachten. 

% Zur besseren Übersicht können Abschnitte in Dateien ausgelagert werden.
\section{Grundlegender Aufbau}
Auf eine Überschrift folgt immer ein Absatz mit Text, \textbf{nie unmittelbar} eine Überschrift niedrigeren Ranges \cite[S.~197]{Balz2008}.\footnote{Hervorhebungen werden durch Kursivschrift oder halbfette Schrift realisiert. Vermieden werden soll Hervorhebung durch Unterstreichen, größere Schrift oder Großbuchstaben \cite[S.~182]{Balz2008}. Ebenfalls ungünstig ist die farbliche Hervorhebung von Wörtern oder Überschriften.}
Daher muss von der Einleitung eine schlüssige Überleitung zum Hauptteil und den anstehenden Unterkapiteln erfolgen.
Falls mehrere Versuche innerhalb des Praktikums durchgeführt werden, sollte an dieser Stelle kurz darauf eingegangen werden.

Ein Hauptkapitel (Kapitel mit nur einer Zählnummer) beginnt jeweils auf einer neuen Seite.
Innerhalb der Gliederung ist zu beachten, dass Gliederungsebene vier nicht überschritten werden sollte.
Die Unterteilung soll möglichst in der dritten Stufe enden \cite[S. 12]{Hering2009}.

Zum Beispiel soll eine Gliederungsebene "`2.2.1.3"' vermieden werden, also wird die Untergliederung bei Ebene "`2.2.1"' beendet.
Desweiteren muss darauf geachtet werden, dass ein Oberpunkt der Gliederung mindestens zwei Unterpunkte besitzt muss \cite[S. 53]{Karm2011}.
Dennoch sollen die Zählnummern innerhalb einer Hierarchiestufe nicht größer als neun werden \cite[S.~12]{Hering2009}.
Die jeweils übergeordneten Kapitel sollen einen Überblick über die Inhalte geben, die in den Unterkapiteln thematisiert sind.
In den Überblicken sollte jedoch kein zusätzliches Wissen vermittelt werden \cite[S.~196]{Balz2008}.

In diesem Dokument wird ein beispielhaftes Inhaltsverzeichnis für einen Praktikumsbericht aufgezeigt, der u.a. im Praktikum Grundlagen der Elektrotechnik an der HSHL angefertigt werden muss.
Die Inhaltsübersicht und das Inhaltsverzeichnis sollen einen Überblick über die Systematik der Arbeit geben und einen "`ersten Eindruck"' vermitteln.
Man sollte sich bewusst sein, dass das Inhaltsverzeichnis die erste Information ist, die ein Leser von der Arbeit aufnimmt \cite[S.~49]{Karm2011}.
Die Gliederung soll bereits als "`roter Faden"' dienen.
Für den Leser ist die Gliederung der wichtigste Schlüssel zum Erfassen des Inhalts.
Wissenschaftliche Arbeiten erfordern ein hohes Maß an Ordnung und Logik, welche sich auch in der Gliederung widerspiegeln muss \cite[S.~9 ~ff.]{Hering2009}.

\noindent
Üblicherweise gliedert sich der inhaltliche Aufbau in Einleitung, Hauptteil und Schluss.
Die Einleitung beinhaltet die Aufgabenstellung und die Motivation der wissenschaftlichen Arbeit (s. Einleitung, S. \pageref{sec:Einleitung}).
Im Hauptteil werden die theoretischen Grundlagen erläutert und der Lösungswegs sowie die Ergebnisse dargestellt.
Die Ergebnisse werden im Schlussteil in Kurzform wiedergegeben und kritisch diskutiert.
Weitere Untergliederungen dieser drei Hauptpunkte können je nach Art der wissenschaftlichen Arbeit und den entsprechenden Anforderungen umgesetzt werden.
In dieser wissenschaftlichen Arbeit werden die benötigten Gliederungspunkte für einen Praktikumsbericht aufgezeigt.

Ein Praktikumsbericht, der die Beschreibung einen Versuchs beinhaltet, sollte unbedingt folgende Unterpunkte beinhalten:
\begin{enumerate}
\item Theoretische Grundlagen 
\item Aufgabenstellung mit Versuchsbeschreibung
\item Versuchsvorbereitung (benötigter Materialien und Beschreibung des Versuchsaufbaus)
\item Versuchsdurchführung
\item Versuchsauswertung (vgl. \cite{Hering2009}, S. 18)
\end{enumerate}

Als allgemeine Gliederung für ingenieurwissenschaftliche Arbeit (z.\,B Bachelor- oder Masterarbeit) empfiehlt sich die Nachfolgende. 

Titelseite\\
Inhaltsverzeichnis
\begin{enumerate}
	\item Einleitung (Einführung in die Aufgabenstellung)
  \item Theoretische Grundlagen
	\item Lösungsweg ggf. Lösungsalternative
  \item Ergebnis
  \item Zusammenfassung (mit anschl. Diskussion der Ergebnisse) und Ausblick
\end{enumerate}
Quellenverzeichnis\\
Anhang\\
Eidesstattliche Erklärung 

%---------------------------------------------
\section{Stilistische Anforderungen}
%---------------------------------------------
Die "`Ich"'- oder "`Wir"'-Perspektive soll im gesamten Textteil vermieden werden.
Stattdessen sollte aus der dritten, neutralen Person oder in Passivkonstruktionen geschrieben werden \cite[S.~130]{Hering2009}.
Die wissenschaftliche Arbeit darf nicht im Sinne eines Erlebnisberichtes sondern als sachorientierte Dokumentation aufgefasst werden. 
Eine Ausnahme besteht im Schlussteil, in dem der Autor eine kritische Auseinandersetzung mit der Thematik einbringen darf.

Ein wesentliches Kriterium für die Verständlichkeit der Inhalte ist ein wissenschaftlich-technischer, konkreter, klarer und präziser Sprachstil.
Es sollte auf keinen Umständen Umgangssprache verwendet werden. 
Folgende Punkte sorgen für bessere Textverständlichkeit:
\begin{itemize}
	\item Kurze Sätze bilden.
	\item Möglichst nur einen Hauptsatz mit einem oder zwei Nebensätzen bzw. zwei mit Semikolon verknüpfte Hauptsätze verwenden. 
	\item Nicht zu viele Nebensätze in Klammern sondern als alleinstehende Sätze oder in Fußnoten unterbringen.
	\item "`Welcher, welche, welches"' als Relativpronomen sind veraltet und sollen ersetzt werden durch "`der, die, das"'. 
	\item Ein- und überleitende Sätze verwenden, um Ihre Leser mit Worten zu führen.
Damit können Sie:\\
\begin{itemize}
	\item die Gliederung bzw. das Inhaltsverzeichnis aufgreifen,
	\item den bisher beschriebenen Sachverhalt zusammenfassen,
	\item zum nächsten Dokumentteil überleiten oder
	\item einen neuen Dokumentteil einleiten.
\end{itemize}
	\item Anschaulich formulieren, in Bildern sprechen!
	Analogien, Metaphern und Vergleiche erzeugen Assoziationen beim Leser.
	Dadurch kann der Leser Ähnlichkeiten und Abweichungen von dem ihm bereits bekannten Wissen leichter erkennen \cite[S.~129]{Hering2009}.
	\item Konjunktive vermeiden.
	\item Apostrophe sparsam im Text verwenden.
	\item Indirekte statt direkte Zitate verwenden.
	\item Wiederholungen von Wörtern in Sätzen vermeiden.
	\item Füllwörter, d.\,h. Wörter die beim Weglassen keinen Einfluss auf den Inhalt haben, vermeiden.
	\item Inhaltliche Blöcke im Text bilden und diese durch Einrückungen voneinander trennen.
\end{itemize}

\noindent
Die Erzählzeit ist das Präsens (Gegenwart).
Das Imperfekt (Vergangenheit) wird nur bei Beschreibung eines früher verwendetes Bauteil, Messverfahren usw., eingesetzt.
Es sollten keine emotionsbeladenen und umgangssprachlichen Formulierungen verwendet werden \cite[S.~131]{Hering2009}.
Zur verständliche Formulierung gehören einleitende Sätze am Anfang eines Dokumentteils.
Zusammenfassende und überleitende Sätze sollten am Ende eines Dokumentteils eingebracht werden.
Diese Sätze können auch Einsatz zur Ein- und Überleitung bei Bildern und Tabellen finden.
Sie beschreiben z.B. welche Informationen und warum als nächstes folgen, wie die folgenden Informationen zu bewerten sind, welchen Stellenwert die zuvor beschriebenen Informationen haben usw. (vgl. \cite{Hering2009}, S. 136 f.).

Technische Geräte sollen grundsätzlich nach ihrer Funktion benannt werden.
Zum Beispiel ist die Bezeichnung Antriebsrad besser als Zahnrad oder Halteplatte eindeutiger als Platte.
Technische Verfahren werden zur Veranschaulichung funktionsgerecht benannt.
So sollte die Bezeichnung "`Variante 1"' in z.\,B. "`elektrisch-mechanische Lösung"' umbenannt werden \cite[S.~131]{Hering2009}.
Dabei muss aber beachtet werden, dass keine Worte erfunden werden dürfen.
Es sollen nur Worte verwendet werden, die im Duden stehen.
Bezeichnungen wie z.B. "`bestromen"' oder "`anflanschen"' sollten in einer wissenschaftlichen Arbeit durch gebräuchliche Ausdrücke ersetzt werden.

Bei der Benutzung von Fachbegriffen, Fremdwörten und Abkürzungen gilt es eine angemessene Verwendung einzuhalten.
Fremdsprachige Wörter sollten vermieden werden.
Nach Möglichkeit sollen die sinngemäßen Bezeichnungen in der deutschen Sprache verwendet werden.
Ausnahme bilden fremdsprachige Fachausdrücke, bei denen eine Übersetzung ins Deutsche übertrieben erscheinen würde (sogenannte Anglizismen vermeiden).
Bei der Verwendung von Fachbegriffen müssen diese bei erstmaliger Verwendung erläutert werden.
Dies kann sowohl innerhalb des Textes als auch in Fußnoten erfolgen.\footnote{In den Fußnoten können Anmerkungen untergebracht werden ohne den Lesefluss zu stören.
In die Fußnoten gehören ausschließlich zusätzliche Informationen und Anmerkungen, die als nennenswert erachtet werden, aber nicht für die Argumentation erforderlich sind. Jede Fußnote beginnt mit einem Großbuchstaben und endet mit einem Punkt \cite[S.~197~f.]{Balz2008}.}
Unbekannte Abkürzungen werden ebenfalls beim ersten Auftreten definiert \cite[S.~136]{Balz2008}.
Sobald eine Abkürzung einmalig eingeführt wurde, muss sie konsequent verwendet werden.
Beim Einführen einer Abkürzung wird zuerst der vollständige Begriff genannt und in Klammern auf die Abkürzung verwiesen wird.
Bei mehr als fünf verwendeten Abkürzungen sollte ein Abkürzungsverzeichnis erstellt werden.


%===================================================================================================================
\subsection{Zitieren}
Sobald auf fremde Literaturquellen referenziert wird, ist die Verwendung einer korrekten Zitierweise und der Aufbau eines Literatur- und Quellenverzeichnis erforderlich.
Die Inhalte fremder Quellen können in einem direkten oder indirekten Zitat wiedergegeben werden.

Bei der wörtlichen Übernahme einer fremden Quelle im eigenen Text wird von einem direkten Zitat gesprochen.
Dieses steht immer in Anführungszeichen und muss buchstaben- und zeichengetreu übernommen werden, auch wenn das Original Rechtschreib- oder sonstige Fehler enthält.

Ein indirektes Zitat liegt vor, wenn fremde Gedanken und Ausführungen sinngemäß übernommen werden.
Es steht nie in Anführungszeichen und wird z.\,B. mit einem "`vgl."' referenziert.

Man spricht von einem Vollbeleg, wenn ein gesondertes Literaturverzeichnis vorliegt und auf die Einträge im Literaturverzeichnis verwiesen wird wie das in dieser wissenschaftlichen Arbeit der Fall ist.
Bei einem Kurzbeleg existiert kein gesondertes Literaturverzeichnis.
Die Quellen der Zitate werden dabei innerhalb des Textes oder in Fußnoten aufgeführt \cite{Kipp2012}.

Sinngemäß zitierte Sachverhalte werden durch die Quellenangabe z.\,B. am Ende eines Teilsatzes, eines Satzes oder Absatzes gekennzeichnet \cite[S.~109~ff.]{Hering2009}.
Eine wissenschaftliche Arbeit stützt sich nicht nur auf Schriftwerke, sondern auch auf andere geschützte Werke der Literatur, Wissenschaft und Kunst.
Der Textteil beinhaltet neben den eigenen und zitierten Inhalten auch Abbildungen, Tabellen und Gleichungen, welche an entsprechenden Stellen in den Text eingebunden werden \cite[S.~196]{Balz2008}.
Die Theorie ist mit entsprechenden Zitaten, Formeln, Tabellen und Bildern aus Fachliteratur zu belegen.
Dabei sollen grundsätzlich Primärquellen den Sekundärquellen vorgezogen werden.
Primärquellen oder Originalquellen sind eigenständige wissenschaftliche Arbeiten, die einen konkreten Betrachtungsgegenstand haben und zum Zeitpunkt der Veröffentlichung neue Erkenntnisse über diesen liefern \cite[S.~81]{Balz2008}.

"`Schwierige und komplexe Inhalte brauchen Wahrnehmungshilfen [...]: Dazu zählen Überschriften, Auszeichnungen, Bilder, Symbole und Strukturbilder, Tabellen und Diagramme, Fotos und multimediale Elemente"' (\cite{Balz2008}, S. 28).
Wenn Formeln und Gleichungen verwendet werden, so werden sie nicht in den Text eingebettet sondern nehmen eine eigene Zeile ein. 
Es ist üblich, die Gleichungen zusätzlich mit Gleichungsnummern am rechten Rand des Satzspiegels zu kennzeichnen, damit man sich im Text auf die jeweilige Gleichung beziehen kann.
Im gesamten Text kann auf die Formeln Bezug genommen werden indem ihre Nummer genannt wird.
Dabei werden die Klammern weggelassen \cite[S.~133~ff.]{Hering2009}.
Nachfolgend wird ein Beispiel für das korrekte Einbringen einer Formel mit dem dazugehörigen Querverweis im Text dargestellt.

Nach Gleichung \ref{eq:Gleichung1} errechnet sich die Gesamtsystemleistung $P_{ges}$ aus dem Produkt von Spannung $U_f$, Strom $I_f$ und der Anzahl der LEDs $n$.

\begin{equation}
P_{ges} = P_E = I_f \cdot U_f \cdot n
\label{eq:Gleichung1}
\end{equation}

\noindent
Formeln treten selten als einzelne Formel wie in Gleichung \ref{eq:Gleichung1} auf.
Meist erscheinen sie im Rahmen umfangreicherer Berechnungen mit mehreren Formeln und Umformungen.
Innerhalb des Textes sollten die Formelgrößen erklärt, Formelumformungen beschrieben und die Ergebnisse aufgegriffen werden.
Die Einflussgrößen, die in einer Formel auftreten, werden im Text kurz vor oder direkt nach der Formel benannt und erläutert \cite[S.~134]{Hering2009}.
Es folgt ein Beispiel für eine umfangreichere Berechnung.

Die LEDs sind mit einem Durchflussstrom von $700\,mA$ und einer Durchflussspannung von $3,31V$ klassifiziert.
Um die benötigte Systemleistung $P_{ges}$ aufzubringen, wird von dem typischen Durchflussstrom von $700\,mA$ ausgegangen.
Bei diesem stellt sich nach Gleichung \ref{eq:Spannung319V} eine Durchflussspannung von $3,19\,V$ ein und erreicht bei der Anzahl $n$ von 14 LEDs somit die Gesamtsystemleistung von $31,25\,W$.
%
\begin{align}
P_{ges} = P_E &= I_f \cdot U_f \cdot n \\
\Leftrightarrow U_f	&= \frac{P_{ges}}{I_f \cdot n}\\
	  &= \frac{31,25W}{700mA \cdot 14}\\
	  &= 3,19V
\label{eq:Spannung319V}
\end{align}

Gleichungen sowie Zahlen mit Einheiten werden kursiv hervorgehoben, falls diese im Textteil erscheinen.
Mit Latex gesetzte Formeln sehen einheitlich aus, da Formelzeichen in der mathematischen Umgebung grundsätzlich kursiv dargestellt werden.

%===================================================================================================================
\subsection{Quellen}
Fertigen Sie Ihr Literaturverzeichnis nach DIN ISO 690:2013-10 ist an.
Die DIN ISO 690:2013-10 sieht im Literaturverzeichnis zwingend die Nennung des Urhebers, Sachtitels, Auflage und Erscheinungsvermerke wie Erscheinungsort, Verlag und Erscheinungsjahr einer Literaturquelle vor.
\newpage
Dies sind Beispiele für Zitate nach DIN ISO 690:2013-10:

\textbf{Bücher}
  \begin{tabbing}
    [1]x \= Mitte  \kill
    $[1]$ \> \textsc{Paul, R.}: \textit{Elektrotechnik 1.} Berlin: Springer-Verlag, 3. Auflage 1993.\\
		\> ISBN: 978-0-387557-53-3\\
		$[2]$ \> \textsc{Balzert, H. u.a.}: \textit{Wissenschaftliches Arbeiten.} \\
		\> Witten: W3L-Verlag, 2. Auflage 2008. ISBN 987-3-937137-59-9\\
		$[3]$ \> \textsc{Eichler, H.-J.; Kronfeldt, H.-D.; Sahm, J.}: \textit{Das neue Physikalische Grund-}\\
		\>\textit{praktikum.} Berlin, Heidelberg: Springer-Verlag, 3. Auflage 2016.\\
		\> ISBN: 978-3-662490-22-8
  \end{tabbing}

\textbf{Internetquellen}
  \begin{tabbing}
    [1]x \= Mitte  \kill
    [4] \> \textsc{Wikipedia}: \textit{Feldlinien.} URL: \url{https://de.wikipedia.org/wiki/Feldlinien}, \\
		\>Abruf: 12.02.2022\\
  \end{tabbing}

\textbf{Zeitschriften}
  \begin{tabbing}
    [1]x \= Mitte  \kill
    [5] \> \textsc{Zhong, X.;Lewis, J. M.}: \textit{A New Method for Autonomous Robot Calibration.} \\
		\>In: Proc. of the IEEE Int. Conf. on Robotics and Automation. \\
		\>IEEE, 1995, S. 1790-1795 Vol.2. DOI: 10.1109/ROBOT.1995.525529.\\
  \end{tabbing}




%================================================================================================================================================
\subsection{Verwendung von Bildern}
Die verbale Darstellung von Informationen ist zwar exakt aber nicht sehr anschaulich.
Nach dem Prinzip "`ein Bild sagt mehr als tausend Worte"' können viele Sachverhalte einfacher und klarer in Form von Bildern dargestellt werden als es durch verbale Beschreibung möglich ist.
Die Inhalte sollten grafisch dargeboten werden, weil sie so leichter aufgenommen und verarbeitet werden können \cite[S.~67]{Hering2009}.
Es ist wichtig eine geeignete Darstellungsform für die Veranschaulichung der theoretischen Inhalte zu finden.
Dies müssen nicht unbedingt Abbildungen, sondern können auch Tabellen, Diagramme oder andere Darstellungsformen sein.

"`Grafiken, Tabellen und Fotos, die aus einer fremden Quelle übernommen wurden, müssen mit einem entsprechenden Verweis auf den Ursprung versehen werden.
Eigenerstellte Inhalte werden nicht explizit gekennzeichnet"' \cite[S.~111]{Balz2008}.
Wie in Abbildung \ref{fig:KompassSensor} zu sehen ist, setzt sich die Abbildungsunterschrift aus Nummerierung, gefolgt von der Benennung und der Quellenangabe zusammen. 
Bei der Abbildungsunterschrift sowie Tabellenüberschriften wird auf das Satzzeichen am Ende des Satzes verzichtet.
In der Regel gilt bei Abbildungen und Tabellen eine horizontal zentrierte Anordnung, da das optische Erscheinungsbild dadurch ausgewogen erscheint \cite[S.~183]{Balz2008}.
Für die Abbildungsunterschrift sowie Tabellenüberschrift gilt ebenfalls eine zentrierte Anordnung.

\begin{figure}[h]
	\centering
		\includegraphics[width=0.5\textwidth]{KompassSensor.jpg}\\
			\caption{Kompasssensor für Lego Mindstorms NXT und EV3 \cite{Lego12}}
		\label{fig:KompassSensor}
\end{figure}

\noindent
Bilder und Tabellen werden entweder durch die ganze wissenschaftliche Ausarbeitung chronologisch durchnummeriert wie es in dieser wissenschaftlichen Arbeit gehandhabt wird.
Oder die Nummern werden aus der Kapitelnummer und einer laufenden Nummer innerhalb des Kapitels kombiniert \cite[S.~15]{Hering2009}.

Jede Abbildung muss als Querverweis im Text referenziert werden \textbf{bevor} diese im Dokument erscheint.
Auf jedes Bild muss im Text mindestens einmal hingewiesen werden.
Dabei soll eine Verbindung zwischen Bild und Text geschaffen werden und das Bild möglichst in der Nähe des Hinweises im Text erscheinen.
Falls von mehreren Stellen auf ein Bild hingewiesen wird, dann sollte das Bild in der Nähe des wichtigsten Hinweises angeordnet werden \cite[S.~76]{Hering2009}.

Bei der Auswahl der Abbildung muss darauf geachtet werden, dass diese in angemessener Druckqualität vorliegt.
Die Qualität kann durch ein Bildbearbeitungsprogramm optimiert werden, indem der Kontrast oder die Schärfe angepasst wird.
Grundsätzlich sollte eine Abbildung nur Informationen beinhalten, die im Text beschrieben werden.
Damit das Bild übersichtlich bleibt, sollte in Abbildungen nur das Wesentliche dargestellt werden.
Durch das Weglassen und Vereinfachen von Elementen in der Abbildung wird beim Leser die Konzentration auf das Wesentliche gelenkt \cite[S.~70~f.]{Hering2009}.

Der Übersichtlichkeit halber sollten auf einer Textseite nicht mehr als zwei Abbildungen, Tabellen oder andere Visualisierungsformen abgebildet werden.
Generell sollten Abbildungen nur dann eingesetzt werden, wenn es notwendig ist, komplexe Inhalte darzustellen, nicht als bloße Dekoration.
Wenn innerhalb der Abbildung schriftliche Anteile enthalten sind, sollte die Schriftgröße leserlich oder optimalerweise so groß wie die verwendete Schriftgröße im Text sein.
Es sollte vermieden werden, dass ein Bild unbekannte Fachbegriffe enthält oder dass in Text und Bild unterschiedliche Bezeichnungen für gleiche Gegenstände oder Sachverhalte auftreten \cite[S. 80]{Hering2009}.




